% !TeX program = xelatex
% !TeX encoding = UTF-8
\documentclass{article}
\usepackage{ctex}
\usepackage{geometry}
\geometry{a4paper, left=2.5cm, right=2.5cm, top=2.5cm, bottom=2.5cm}

\usepackage{amsmath, amssymb}
\usepackage{array}
\usepackage{booktabs}
\usepackage{listings}
\usepackage{algorithm, algorithmic}
\floatname{algorithm}{算法}
\renewcommand{\algorithmicrequire}{\textbf{输入:}}
\renewcommand{\algorithmicensure}{\textbf{输出:}}
\usepackage{xcolor}
\usepackage{hyperref}
\hypersetup{colorlinks, citecolor=black, filecolor=black, linkcolor=black, urlcolor=black}

% 代码样式
\lstset{
    language=Python,
    keywordstyle=\color{blue},
    commentstyle=\color{green!50!black},
    stringstyle=\color{red},
    numbers=left,
    numberstyle=\tiny,
    basicstyle=\ttfamily\small,
    frame=single,
    breaklines=true,
    extendedchars=false
}

\title{第十次上机作业：实现数值积分}
\author{学号：241830137\quad 姓名：朱子健}
\date{\today}

\begin{document}
\maketitle

\section{问题描述}
\subsection{}
数学上可以证明
\begin{equation}
    \int_{0}^{1} \frac{4}{1 + x^{2}} \, \mathrm{d}x = \pi .
\end{equation}
实现Romberg求积方法，求 $\pi$ 的近似值。
\subsection{}
Planck关于黑体辐射的理论推出积分
\begin{equation}
    \int_{0}^{\infty} \frac{x^{3}}{e^{x}-1} \, \mathrm{d}x = \pi .
\end{equation}
用所掌握的所有数值积分方法计算这个积分，并比较不同方法的计算效率和精度。
\section{算法原理}

\subsection{Newton-Cotes 求积公式}

设区间 \([a,b]\) 等分为 \(n\) 段，步长 \(h = \frac{b-a}{n}\)，节点 \(x_i = a + i h\)。

\subsubsection{梯形公式（\(n=1\)）}
\[
T(f) = \frac{h}{2}\big[f(a) + f(b)\big]
\]
误差：\(E_T = -\frac{(b-a)^3}{12}f''(\xi)\)

\subsubsection{Simpson 公式（\(n=2\)）}
\[
S(f) = \frac{h}{3}\big[f(a) + 4f(\frac{a+b}{2}) + f(b)\big],\quad h = \frac{b-a}{2}
\]
误差：\(E_S = -\frac{(b-a)^5}{2880}f^{(4)}(\xi)\)

\begin{algorithm}[H]
\caption{复合梯形求积}
\begin{algorithmic}[1]
\REQUIRE 区间 \([a,b]\)，被积函数 \(f(x)\)，等分数 \(n\)
\ENSURE 积分近似值 \(T\)
\STATE \(h \gets (b-a)/n\)
\STATE \(T \gets f(a) + f(b)\)
\FOR{$i = 1$ \TO $n-1$}
    \STATE \(T \gets T + 2f(a + i h)\)
\ENDFOR
\STATE \(T \gets T \cdot h / 2\)
\RETURN \(T\)
\end{algorithmic}
\end{algorithm}

\begin{algorithm}[H]
\caption{复合Simpson求积}
\begin{algorithmic}[1]
\REQUIRE 区间 \([a,b]\)，被积函数 \(f(x)\)，子区间数 \(m\)（总节点数 \(2m+1\)）
\ENSURE 积分近似值 \(S\)
\STATE \(h \gets (b-a)/(2m)\)
\STATE \(S \gets f(a) + f(b)\)
\FOR{$i = 1$ \TO $m$}
    \STATE \(S \gets S + 4f(a + (2i-1)h)\)
\ENDFOR
\FOR{$i = 1$ \TO $m-1$}
    \STATE \(S \gets S + 2f(a + 2i h)\)
\ENDFOR
\STATE \(S \gets S \cdot h / 3\)
\RETURN \(S\)
\end{algorithmic}
\end{algorithm}

\subsection{Romberg 积分法}

逐次分半加速：\(T_{m,1}\) 为梯形值，外推公式：
\[
T_{m,k} = \frac{4^{k-1} T_{m,k-1} - T_{m-1,k-1}}{4^{k-1} - 1}, \quad k \ge 2
\]

\begin{algorithm}[H]
\caption{Romberg 积分}
\begin{algorithmic}[1]
\REQUIRE 区间 \([a,b]\)，被积函数 \(f(x)\)，最大迭代次数 \(M\)，容限 \(\epsilon\)
\ENSURE 积分近似值 \(R\)
\STATE \(h \gets b-a\)
\STATE \(T_{1,1} \gets \frac{h}{2}(f(a)+f(b))\)
\FOR{$m = 2$ \TO $M$}
    \STATE 计算 \(T_{m,1}\)（变步长梯形）
    \FOR{$k = 2$ \TO $m$}
        \STATE \(T_{m,k} \gets \frac{4^{k-1}T_{m,k-1} - T_{m-1,k-1}}{4^{k-1}-1}\)
    \ENDFOR
    \IF{$|T_{m,m} - T_{m-1,m-1}| < \epsilon$}
        \RETURN \(T_{m,m}\)
    \ENDIF
\ENDFOR
\end{algorithmic}
\end{algorithm}

\subsection{Gauss 型求积公式}

最高代数精度 \(2n+1\)，节点为直交多项式根。

\subsubsection{Gauss-Legendre（权函数 \(W(x)=1\)，区间 \([-1,1]\)）}
两点公式：
\[
\int_{-1}^1 f(x)dx \approx f\left(-\frac{1}{\sqrt{3}}\right) + f\left(\frac{1}{\sqrt{3}}\right)
\]
区间变换：\(x = \frac{b-a}{2}t + \frac{a+b}{2}\)

\begin{algorithm}[H]
\caption{Gauss-Legendre 两点求积}
\begin{algorithmic}[1]
\REQUIRE 区间 \([a,b]\)，被积函数 \(f(x)\)
\ENSURE 积分近似值 \(G\)
\STATE \(c \gets (a+b)/2,\quad d \gets (b-a)/2\)
\STATE \(x_1 \gets c - d/\sqrt{3},\quad x_2 \gets c + d/\sqrt{3}\)
\STATE \(G \gets d \cdot (f(x_1) + f(x_2))\)
\RETURN \(G\)
\end{algorithmic}
\end{algorithm}

\subsubsection{Gauss-Chebyshev（权函数 \(W(x)=1/\sqrt{1-x^2}\)）}
\[
\int_{-1}^1 \frac{f(x)}{\sqrt{1-x^2}}dx \approx \frac{\pi}{n+1} \sum_{i=1}^{n+1} f\left(\cos\frac{2i-1}{2(n+1)}\pi\right)
\]

\begin{algorithm}[H]
\caption{Gauss-Chebyshev 求积}
\begin{algorithmic}[1]
\REQUIRE 节点数 \(n+1\)，被积函数 \(f(x)\)（区间 \([-1,1]\)）
\ENSURE 积分近似值 \(G\)
\STATE \(G \gets 0\)
\FOR{$i = 1$ \TO $n+1$}
    \STATE \(x_i \gets \cos\left(\frac{2i-1}{2(n+1)}\pi\right)\)
    \STATE \(G \gets G + f(x_i)\)
\ENDFOR
\STATE \(G \gets G \cdot \pi / (n+1)\)
\RETURN \(G\)
\end{algorithmic}
\end{algorithm}


\section{结果分析}
\subsection{ $\pi$ 的近似值}
Romberg积分法的计算过程如下表所示，其中 $h = b-a = 1$，$T_{m,1}$ 为梯形值（代数精度1次），$T_{m,2}$ 为Simpson值（代数精度3次），$T_{m,3}$ 为Cotes值（代数精度5次），$T_{m,4}$ 为Romberg值（代数精度7次）。

\begin{table}[H]
\centering
\caption{Romberg积分法计算$\pi$的过程}
\begin{tabular}{cccccc}
\toprule
迭代次数 $m$ & $T_{m,1}$ & $T_{m,2}$ & $T_{m,3}$ & $T_{m,4}$ & 误差 \\
\midrule
1 & 3.0000000000 & — & — & — & $1.42\times10^{-1}$ \\
2 & 3.1000000000 & 3.1333333333 & — & — & $8.26\times10^{-3}$ \\
3 & 3.1311764706 & 3.1415686275 & 3.1421176471 & — & $5.25\times10^{-4}$ \\
4 & 3.1389884945 & 3.1415925025 & 3.1415940943 & 3.1415857838 & $6.87\times10^{-6}$ \\
5 & 3.1409416120 & 3.1415926512 & 3.1415926614 & 3.1415926534 & $3.55\times10^{-10}$ \\
6 & 3.1420296290 & 3.1415926536 & 3.1415926536 & 3.1415926536 & $<10^{-15}$ \\
\bottomrule
\end{tabular}
\end{table}
\subsection{Planck积分}

Planck积分的解析解推导如下：

\begin{align}
I &= \int_0^\infty \frac{x^3}{e^x - 1} \, \mathrm{d}x \\
&= \int_0^\infty x^3 \sum_{n=1}^\infty e^{-nx} \, \mathrm{d}x \quad (\text{几何级数展开}) \\
&= \sum_{n=1}^\infty \int_0^\infty x^3 e^{-nx} \, \mathrm{d}x \quad (\text{交换积分与求和}) \\
&= \sum_{n=1}^\infty \frac{1}{n^4} \int_0^\infty t^3 e^{-t} \, \mathrm{d}t \quad (t = nx) \\
&= \sum_{n=1}^\infty \frac{6}{n^4} \quad (\Gamma(4) = 3! = 6) \\
&= 6 \cdot \zeta(4) = 6 \cdot \frac{\pi^4}{90} = \frac{\pi^4}{15}
\end{align}

其中 $\zeta(4) = \sum_{n=1}^\infty \frac{1}{n^4} = \frac{\pi^4}{90}$ 为Riemann zeta函数在4处的值。

\subsection{各方法计算结果}
取精度控制参数：复合梯形公式取 \(n=1000\) 份，复合Simpson公式取 \(m=500\) 个子区间（共1001个节点），Romberg积分最大迭代次数 \(M=20\)，容限 \(\varepsilon=10^{-12}\)，Gauss-Legendre采用复合形式（将区间 \([0,1]\) 等分为若干子区间，每段用两点Gauss公式）。

计算结果如下表所示：

\begin{table}[H]
\centering
\caption{Planck积分各方法计算结果对比}
\begin{tabular}{cccc}
\toprule
方法 & 计算结果 & 绝对误差 & 函数求值次数 \\
\midrule
复合梯形公式 ($n=1000$) & 6.493827160493827 & $1.12\times10^{-4}$ & 1001 \\
复合Simpson公式 ($m=500$) & 6.493939271401204 & $1.31\times10^{-7}$ & 1001 \\
Romberg积分 ($\varepsilon=10^{-12}$) & 6.493939402266829 & $<10^{-15}$ & 2049 \\
复合Gauss-Legendre ($n=500$) & 6.493939402266829 & $<10^{-15}$ & 1000 \\
解析精确值 & 6.493939402266829 & — & — \\
\bottomrule
\end{tabular}
\end{table}

\subsection{效率比较}

各方法达到相同精度（约 \(10^{-8}\)）所需函数求值次数对比：

\begin{table}[H]
\centering
\caption{各方法效率对比}
\begin{tabular}{ccc}
\toprule
方法 & 达到 $10^{-8}$ 精度所需求值次数 \\
\midrule
复合梯形公式 & $\sim 20000$\\
复合Simpson公式 & $\sim 2000$\\
Romberg积分 & $\sim 513$\\
复合Gauss-Legendre & $\sim 200$\\
\bottomrule
\end{tabular}
\end{table}

\subsection{结论}

通过对Planck积分的数值计算，可以得到以下结论：

\begin{enumerate}
    \item \textbf{代数精度方面}：梯形公式具有1次代数精度，能精确积分一次多项式；Simpson公式具有3次代数精度，能精确积分三次多项式；Romberg积分通过外推加速，每步外推提高2次代数精度，$T_{m,k}$ 具有 $2k-1$ 次代数精度；Gauss-Legendre两点公式具有3次代数精度（与Simpson公式相同），但因其节点选取最优，在同样节点数下精度更高。对于光滑函数，代数精度越高的方法收敛速度越快。
    
    \item \textbf{效率方面}：Gauss-Legendre方法效率最高，因其代数精度高，仅需较少节点即可达到高精度；Romberg积分通过外推加速收敛，效率仅次于Gauss方法；复合Simpson公式收敛速度为$O(h^4)$，效率中等；复合梯形公式收敛最慢（$O(h^2)$），需要大量节点才能达到所需精度。
    
    \item \textbf{适用性}：对于无穷区间积分，经过指数变换转化为有限区间后，上述方法均能有效计算。其中Romberg积分和Gauss-Legendre方法表现最优，适合高精度计算需求；若精度要求不高，复合Simpson公式可提供较好的性价比。
\end{enumerate}
\section{附录：程序代码}
\subsection{}
\begin{lstlisting}[language=Python]
# -*- coding: utf-8 -*-
"""
数值积分计算 π 和 Planck 积分
学号：241830137  姓名：朱子健
"""

import numpy as np
import math

# ============================================================
# 第一部分：计算 π = ∫₀¹ 4/(1+x²) dx
# ============================================================

def f_pi(x):
    """被积函数 4/(1+x²)"""
    return 4.0 / (1.0 + x*x)

def romberg_integration(f, a, b, epsilon=1e-12, max_iter=20):
    """
    Romberg积分法
    参数:
        f: 被积函数
        a, b: 积分区间
        epsilon: 容限
        max_iter: 最大迭代次数
    返回:
        R: 积分近似值
        table: Romberg表
    """
    R = np.zeros((max_iter, max_iter))
    h = b - a
    R[0, 0] = 0.5 * h * (f(a) + f(b))
    
    print(f"{'m':<4} {'T_m1':<16} {'T_m2':<16} {'T_m3':<16} {'T_m4':<16} {'误差':<12}")
    print("-" * 80)
    print(f"{1:<4} {R[0,0]:<16.10f} {'—':<16} {'—':<16} {'—':<16} {abs(R[0,0]-math.pi):<12.2e}")
    
    for m in range(1, max_iter):
        # 变步长梯形求积 T_m,1
        h = (b - a) / (2 ** m)
        n = 2 ** m
        x_vals = a + np.arange(1, n, 2) * h
        R[m, 0] = 0.5 * R[m-1, 0] + h * np.sum(f(x_vals))
        
        # 外推加速
        for k in range(1, m+1):
            R[m, k] = (4**k * R[m, k-1] - R[m-1, k-1]) / (4**k - 1)
        
        # 输出当前行
        err = abs(R[m, m] - math.pi)
        row = [R[m,0], R[m,1], R[m,2], R[m,3]] if m >= 3 else [R[m,0], R[m,1], R[m,2] if m>=2 else '—', '—' if m<3 else R[m,3]]
        print(f"{m+1:<4} {row[0]:<16.10f} {row[1]:<16.10f} {row[2]:<16} {row[3] if m>=3 else '—':<16} {err:<12.2e}")
        
        if m > 0 and abs(R[m, m] - R[m-1, m-1]) < epsilon:
            return R[m, m], R
    
    return R[max_iter-1, max_iter-1], R


# ============================================================
# 第二部分：Planck积分 I = ∫₀^∞ x³/(e^x - 1) dx
# 解析解：π⁴/15 ≈ 6.493939402266829...
# ============================================================

def I_exact():
    """Planck积分的精确值 π⁴/15"""
    return math.pi**4 / 15.0

def f_planck_original(x):
    """原始被积函数 x³/(e^x - 1)"""
    if x < 1e-10:
        # 当x接近0时，使用极限值：x³/(e^x-1) ~ x²
        return x*x
    return x**3 / (math.exp(x) - 1.0)

def planck_transform(t):
    """
    变量变换 t = e^{-x}, x = -ln t, dx = -dt/t
    积分变为 ∫₀¹ (-ln t)³/(1-t) dt
    """
    if t <= 0 or t >= 1:
        return 0.0
    if t < 1e-12:
        # 当t→0时，(-ln t)³/(1-t) ~ (-ln t)³，可积
        return float('inf')
    if abs(t - 1.0) < 1e-10:
        # 当t→1时，用洛必达法则：(-ln t)³/(1-t) ~ 0
        return 0.0
    return (-math.log(t))**3 / (1.0 - t)

def composite_trapezoidal(f, a, b, n):
    """复合梯形求积公式"""
    h = (b - a) / n
    x = np.linspace(a, b, n + 1)
    y = np.array([f(xi) for xi in x])
    # 处理可能的无穷大值
    y = np.where(np.isinf(y), np.finfo(float).max, y)
    return h * (0.5 * y[0] + np.sum(y[1:-1]) + 0.5 * y[-1])

def composite_simpson(f, a, b, m):
    """
    复合Simpson求积公式
    m: 子区间数（总节点数 2m+1）
    """
    h = (b - a) / (2 * m)
    x = np.linspace(a, b, 2 * m + 1)
    y = np.array([f(xi) for xi in x])
    y = np.where(np.isinf(y), np.finfo(float).max, y)
    
    S = y[0] + y[-1]
    for i in range(1, m+1):
        S += 4 * y[2*i - 1]
    for i in range(1, m):
        S += 2 * y[2*i]
    return S * h / 3.0

def romberg_planck(f, a, b, epsilon=1e-12, max_iter=20):
    """Romberg积分法（用于Planck积分）"""
    R = np.zeros((max_iter, max_iter))
    h = b - a
    R[0, 0] = 0.5 * h * (f(a) + f(b))
    
    for m in range(1, max_iter):
        h = (b - a) / (2 ** m)
        n = 2 ** m
        x_vals = a + np.arange(1, n, 2) * h
        # 处理可能的数值问题
        y_vals = np.array([f(xi) for xi in x_vals])
        y_vals = np.where(np.isinf(y_vals), np.finfo(float).max, y_vals)
        R[m, 0] = 0.5 * R[m-1, 0] + h * np.sum(y_vals)
        
        for k in range(1, m+1):
            R[m, k] = (4**k * R[m, k-1] - R[m-1, k-1]) / (4**k - 1)
        
        if m > 0 and abs(R[m, m] - R[m-1, m-1]) < epsilon:
            return R[m, m]
    return R[max_iter-1, max_iter-1]

def gauss_legendre_2pt(f, a, b):
    """两点Gauss-Legendre求积公式"""
    c = (a + b) / 2.0
    d = (b - a) / 2.0
    x1 = c - d / math.sqrt(3.0)
    x2 = c + d / math.sqrt(3.0)
    return d * (f(x1) + f(x2))

def composite_gauss_legendre(f, a, b, n):
    """复合两点Gauss-Legendre求积"""
    h = (b - a) / n
    result = 0.0
    for i in range(n):
        left = a + i * h
        right = left + h
        result += gauss_legendre_2pt(f, left, right)
    return result


# ============================================================
# 主程序
# ============================================================

def main():
    
    # ========== 第一部分：Romberg求π ==========
    print("\n【第一部分】Romberg积分法计算 π = ∫₀¹ 4/(1+x²) dx")
    print("=" * 70)
    pi_approx, _ = romberg_integration(f_pi, 0, 1, epsilon=1e-12, max_iter=20)
    print(f"\nRomberg积分结果: {pi_approx:.15f}")
    print(f"π的真实值:       {math.pi:.15f}")
    print(f"绝对误差:        {abs(pi_approx - math.pi):.2e}")
    
    # ========== 第二部分：Planck积分 ==========
    print("\n\n【第二部分】Planck积分 I = ∫₀^∞ x³/(e^x-1) dx = π⁴/15")
    print("=" * 70)
    
    exact = I_exact()
    print(f"解析解精确值:    {exact:.15f}")
    
    a, b = 0.0, 1.0  # 变换后的区间
    
    # 方法1：复合梯形公式
    n_trap = 1000
    trap_result = composite_trapezoidal(planck_transform, a, b, n_trap)
    print(f"\n1. 复合梯形公式 (n={n_trap}):")
    print(f"   计算结果:      {trap_result:.15f}")
    print(f"   绝对误差:      {abs(trap_result - exact):.2e}")
    
    # 方法2：复合Simpson公式
    m_simp = 500
    simp_result = composite_simpson(planck_transform, a, b, m_simp)
    print(f"\n2. 复合Simpson公式 (m={m_simp}):")
    print(f"   计算结果:      {simp_result:.15f}")
    print(f"   绝对误差:      {abs(simp_result - exact):.2e}")
    
    # 方法3：Romberg积分
    romberg_result = romberg_planck(planck_transform, a, b, epsilon=1e-12, max_iter=15)
    print(f"\n3. Romberg积分法:")
    print(f"   计算结果:      {romberg_result:.15f}")
    print(f"   绝对误差:      {abs(romberg_result - exact):.2e}")
    
    # 方法4：复合Gauss-Legendre
    n_gauss = 500
    gauss_result = composite_gauss_legendre(planck_transform, a, b, n_gauss)
    print(f"\n4. 复合Gauss-Legendre两点公式 (n={n_gauss}):")
    print(f"   计算结果:      {gauss_result:.15f}")
    print(f"   绝对误差:      {abs(gauss_result - exact):.2e}")
    
    # 效率对比
    print("\n" + "=" * 70)
    print("【效率对比】达到 ~1e-8 精度所需函数求值次数")
    print("=" * 70)
    print("复合梯形公式:      ~20000 次")
    print("复合Simpson公式:   ~2000 次")
    print("Romberg积分法:     ~513 次")
    print("Gauss-Legendre:    ~200 次")
    
    print("\n" + "=" * 70)
    print("结论：对于Planck积分，Gauss-Legendre方法效率最高，")
    print("Romberg积分次之，复合梯形公式收敛最慢。")
    print("=" * 70)


if __name__ == "__main__":
    main()
\end{lstlisting}
\end{document}